\documentclass[11pt]{article}
\usepackage[a4paper,margin=1.9cm]{geometry}
\usepackage{fontspec}
\setmainfont{Latin Modern Roman}
\usepackage{titlesec}
\usepackage{enumitem}
\usepackage[table]{xcolor}
\usepackage{longtable}
\usepackage{booktabs}
\usepackage{array}
\usepackage{pifont}
\usepackage{amssymb}
\usepackage{tikz}
\usepackage{tcolorbox}
\usepackage{hyperref}
\usepackage{fancyhdr}
\usepackage{parskip}
\usepackage{rotating}

\definecolor{gxteal}{HTML}{0F766E}
\definecolor{gxcyan}{HTML}{0284C7}
\definecolor{gxslate}{HTML}{334155}
\definecolor{gxamber}{HTML}{B45309}
\definecolor{gxred}{HTML}{B91C1C}
\definecolor{gxgreen}{HTML}{15803D}
\definecolor{gxviolet}{HTML}{6D28D9}

\hypersetup{colorlinks=true, linkcolor=gxteal, urlcolor=gxcyan,
  pdftitle={GlobeX SIH 2026 Final Problem Statement Report}, pdfauthor={GlobeX Team}}

\titleformat{\section}{\large\bfseries\color{gxteal}}{\thesection}{0.6em}{}
\titleformat{\subsection}{\normalsize\bfseries\color{gxslate}}{\thesubsection}{0.6em}{}
\titlespacing*{\section}{0pt}{1.7em}{0.7em}

\pagestyle{fancy}
\fancyhf{}
\renewcommand{\headrulewidth}{0.4pt}
\fancyhead[L]{\small\color{gxslate}GlobeX --- SIH 2026 Problem Statement Report}
\fancyhead[R]{\small\color{gxslate}Final}
\fancyfoot[C]{\small\thepage}

% ---- tick symbols (visual, no text) ----
\newcommand{\yes}{{\color{gxgreen}\ding{108}}}      % full use
\newcommand{\pmid}{{\color{gxcyan}\ding{109}}}      % partial use
\newcommand{\non}{{\color{black!22}\ding{109}}}     % unused
\newcommand{\bld}{{\color{gxamber}\ding{53}}}       % must be built

\newtcolorbox{callout}[1]{colback=gxteal!5, colframe=gxteal!55, boxrule=0.6pt, arc=2pt,
  left=6pt, right=6pt, top=5pt, bottom=5pt, fonttitle=\bfseries\small, coltitle=white,
  colbacktitle=gxteal, title={#1}}
\newtcolorbox{uspbox}[1]{colback=gxviolet!5, colframe=gxviolet!55, boxrule=0.6pt, arc=2pt,
  left=6pt, right=6pt, top=5pt, bottom=5pt, fonttitle=\bfseries\small, coltitle=white,
  colbacktitle=gxviolet, title={#1}}
\newtcolorbox{warnbox}[1]{colback=gxamber!6, colframe=gxamber!65, boxrule=0.6pt, arc=2pt,
  left=6pt, right=6pt, top=5pt, bottom=5pt, fonttitle=\bfseries\small, coltitle=white,
  colbacktitle=gxamber, title={#1}}

\newcommand{\psbanner}[6]{%
  \vspace{0.2em}
  {\footnotesize\color{gxslate}\textbf{Org:} #1 \ $\vert$ \ \textbf{Theme:} #2 \ $\vert$ \ \textbf{Category:} #3 \ $\vert$ \ \textbf{Similarity rank:} #4 \ $\vert$ \ \textbf{Coverage:} #5 \ $\vert$ \ \textbf{Deadline:} #6}
  \par\vspace{0.6em}
}

\begin{document}

% ================= TITLE =================
\begin{titlepage}
\centering
\vspace*{1.6cm}
{\Huge\bfseries\color{gxteal} SIH 2026 --- Final Problem}\\[0.4em]
{\Huge\bfseries\color{gxteal} Statement Report for GlobeX}\\[1.8em]
{\Large\color{gxslate} Six candidates, measured against what GlobeX actually is}\\[2.6em]
{\large\color{gxslate} Why each one fits \ $\vert$ \ How we would build it \ $\vert$ \ What makes us win}\\[3em]
{\normalsize\color{gxslate} 26132 \ \ 26100 \ \ 26102 \ \ 26021 \ \ 26002 \ \ 26090}\\[2.4em]
{\normalsize\color{gxslate} Prepared: 1 September 2026}\\[0.3em]
{\normalsize\color{gxslate} All 229 published problem statements searched}
\vfill
{\small\color{gxslate} Candidate selection was not done by keyword. A profile of GlobeX's verified
architecture was scored against all 229 problem statements using TF-IDF cosine similarity --- the same
retrieval method GlobeX's own RAG engine uses --- and the ranking is reported in Section~2.\\
One problem statement came out a clear outlier at the top.}
\end{titlepage}

% ================= SUMMARY =================
\section{The short version}

\begin{callout}{What the search found}
\textbf{PS 26132 (Strengthening market linkages and price discovery for farmers)} ranked \textbf{first out
of all 229} problem statements, and not narrowly --- its similarity score is \textbf{65\% higher than the
second-placed statement}. Reading it confirms why: it asks for market intelligence, price forecasting,
verified buyer matching, lot creation, logistics coordination, payment tracking, dispute handling and
transparent transaction records. That is GlobeX's pipeline, moved from cross-border trade to domestic
agriculture. It is the closest thing to "GlobeX, restated as a problem statement" in the entire dataset.
\end{callout}

\vspace{0.4em}
\renewcommand{\arraystretch}{1.45}
\noindent\begin{longtable}{@{}p{1.15cm}p{4.6cm}p{1.5cm}p{1.7cm}p{6.6cm}@{}}
\toprule
\textbf{PS} & \textbf{Short title} & \textbf{Sim.\ rank} & \textbf{Coverage} & \textbf{The one-line case} \\
\midrule
\endhead
\rowcolor{gxteal!8}
26132 & Market linkage \& price discovery for farmers & \textbf{1 / 229} & 9 of 9 & Our product, restated for agriculture. Escrow answers a problem the PS names out loud. \\
26100 & GeM bid compliance verification & 4 / 229 & 14 of 14 & Uses the most of GlobeX at once; the PS itself demands an audit trail. \\
26102 & MPLAD anomaly \& fraud detection & 31 / 229 & 9 of 9 & Our strongest trained model, transferred almost unchanged. \\
26021 & Honey Chain traceability & 30 / 229 & 7 of 8 & Blockchain provenance at the centre, with escrow as the twist. \\
26002 & NER smart logistics platform & 40 / 229 & 6 of 8 & Strong analytics fit; live GPS is new build. \\
26125 & Blockchain identity, access \& asset platform & \textbf{1 / 229}\newline{\scriptsize(chain profile)} & 6 of 8 & The pure blockchain twin --- deepest chain fit in the dataset, but zero ML content. \\
\rowcolor{gxcyan!7}
26130 & Industrial approvals \& compliance single window & 9 / 229\newline{\scriptsize(compliance)} & 8 of 8 & \textbf{New find.} 26100's mirror image, from the applicant's side. \\
\rowcolor{gxcyan!7}
26045 & IP-SAKTI: multilingual source-cited RAG assistant & 4 / 229\newline{\scriptsize(compliance)} & 6 of 7 & \textbf{New find.} The closest twin to our citation-enforced RAG engine. \\
26090 & Artisan market linkage app & 15 / 229 & 6 of 9 & Marketplace core fits; generative AI and mobile are new. \\
\bottomrule
\end{longtable}

% ================= SIMILARITY =================
\section{Finding GlobeX's twin among 229 problem statements}

A profile of GlobeX's verified capabilities --- market discovery, demand and price forecasting,
multi-criteria ranking, counterparty matching and trust scoring, risk and anomaly detection, compliance
retrieval, document verification, blockchain evidence, escrow with disputes, logistics estimation ---
was vectorised alongside all 229 problem statements and scored by cosine similarity.

\vspace{0.5em}
\noindent\begin{minipage}[t]{0.55\textwidth}
\renewcommand{\arraystretch}{1.25}
{\footnotesize
\begin{tabular}{@{}p{0.6cm}p{1.15cm}p{1.1cm}p{4.6cm}@{}}
\toprule
\textbf{\#} & \textbf{PS} & \textbf{Score} & \textbf{Title (short)} \\
\midrule
\rowcolor{gxteal!12}
1 & 26132 & \textbf{0.076} & Farmer market linkage \& price discovery \\
2 & 26188 & 0.046 & Fake identity \& document screening \\
3 & 26006 & 0.042 & Freight forecasting, vessel chartering \\
4 & 26100 & 0.041 & GeM bid compliance verification \\
5 & 26089 & 0.037 & Cooperative gig services platform \\
6 & 26229 & 0.034 & Kabadiwala Connect recycling chain \\
7 & 26182 & 0.033 & Crypto wallet attribution \\
8 & 26228 & 0.033 & CV integrity assurance \\
9 & 26033 & 0.032 & Farmer/consumer marketplace \\
10 & 26018 & 0.032 & Land record digitisation \\
11 & 26108 & 0.032 & Applicable-standards recommender \\
12 & 26099 & 0.031 & Material code harmonisation \\
13 & 26046 & 0.031 & Clinical trials dashboard \\
14 & 26024 & 0.031 & Coal mine governance \\
15 & 26090 & 0.030 & Artisan market linkage app \\
\bottomrule
\end{tabular}}
\end{minipage}\hfill
\begin{minipage}[t]{0.42\textwidth}
\vspace{0pt}
\begin{callout}{Reading the numbers}
Absolute cosine values are low for everything --- a short capability profile against long
problem-statement prose produces sparse overlap. What matters is the \emph{ordering} and the
\emph{gap}, and the gap at the top is decisive: 26132 sits 65\% above second place, while ranks 2--15
are packed within a narrow band of each other.

\vspace{0.5em}
Note also rank 9: PS 26033 is the thin, three-line version of the same farmer-marketplace idea that
26132 states properly. Where two statements describe the same problem, take the one with the
detail --- the rubric rewards root-cause understanding, and 26132 gives us material to work with.
\end{callout}
\end{minipage}

\vspace{0.8em}
\noindent\textbf{Verdict on "is there a project like GlobeX in here?"} Yes --- one. PS 26132 is a genuine
product-level twin. Everything below it matches on \emph{technology} rather than on \emph{product shape}.

\subsection*{Extended sweep: four profiles, not one}
A single product-shaped query favours marketplace problem statements and can miss a statement that is a
deep match to one \emph{part} of GlobeX. The search was therefore re-run with four separate capability
profiles across all 229 statements. Each column below is an independent ranking.

\vspace{0.6em}
\renewcommand{\arraystretch}{1.3}
\noindent{\footnotesize
\begin{tabular}{@{}p{0.5cm}p{3.55cm}p{3.55cm}p{3.55cm}p{3.55cm}@{}}
\toprule
\textbf{\#} & \textbf{Product shape} & \textbf{Blockchain / trust} & \textbf{ML / anomaly / risk} & \textbf{Compliance / RAG} \\
\midrule
1 & \cellcolor{gxteal!14}\textbf{26132} farmer market linkage & \cellcolor{gxviolet!14}\textbf{26125} blockchain identity \& assets & 26170 component burn-in anomaly (ISRO) & \cellcolor{gxteal!10}\textbf{26100} GeM bid compliance \\
2 & 26033 farmer marketplace (thin) & 26194 student innovation (open brief) & 26070 AI/ML general & 26107 BIS standards assistant \\
3 & 26089 cooperative gig platform & 26211 student innovation (hardware) & 26071 rainfall early warning & 26188 fake identity screening \\
4 & 26229 Kabadiwala recycling & \textbf{26021} Honey Chain & 26017 land acquisition delay & \cellcolor{gxcyan!14}\textbf{26045} IP-SAKTI RAG assistant \\
5 & 26006 freight forecasting & 26141 digital signature security & \textbf{26103} project monitoring & \textbf{26024} coal mine governance \\
6 & \textbf{26090} artisan linkage & 26018 land record digitisation & 26178 environmental monitoring & \textbf{26190} legal document management \\
7 & \textbf{26136} startup procurement & 26228 CV integrity assurance & 26025 mine subsidence & 26018 land record digitisation \\
8 & 26028 railway dynamic ETA & 26046 clinical trials dashboard & 26073 weather station anomaly & 26035 test report generation \\
9 & 26164 cryptographic discovery & \textbf{26190} legal document management & \textbf{26102} MPLAD anomaly & \cellcolor{gxcyan!14}\textbf{26130} industrial approvals \\
10 & 26146 bitcoin traffic analysis & \textbf{26100} GeM bid compliance & 26183 crypto fraud exchanges & 26036 metrology verification \\
\bottomrule
\end{tabular}}

\vspace{0.9em}
\noindent\textbf{What the extended sweep changed.} Three things.

\vspace{0.4em}
\noindent\textbf{1. PS 26125 is the blockchain twin, and it is not close.} It ranks \textbf{first of 229}
on the chain/trust profile with a score roughly \emph{70\% above} second place --- the same kind of
decisive gap 26132 shows on product shape. Between them, the two statements bracket GlobeX: 26132 matches
what the product \emph{is}, 26125 matches what the chain layer \emph{does}. Note the mirror-image
weakness though: 26125 ranks \textbf{200th of 229} on the ML profile, because it contains essentially no
machine learning at all.

\vspace{0.4em}
\noindent\textbf{2. Two genuinely new candidates surfaced} --- PS 26130 and PS 26045, both invisible to
the original product-shaped query and both strong enough to earn a place in the comparison (Section~10).

\vspace{0.4em}
\noindent\textbf{3. Nothing better was found on the ML axis.} The ML profile's top hits are real anomaly
problems, but they sit in domains we have no data for and no product story in --- component burn-in
testing at ISRO, rainfall inundation, weather-station telemetry, mine subsidence. PS 26102 remains the
best ML pick because it combines the same modelling shape with published data and a defensible
blockchain angle.

% ================= TICK MATRIX =================
\section{Which parts of GlobeX each problem statement actually uses}

\vspace{-0.3em}
\noindent{\footnotesize
\yes\ = core to the solution \quad
\pmid\ = partially used \quad
\non\ = not used \quad
\bld\ = must be built new}

\vspace{0.8em}
\renewcommand{\arraystretch}{1.6}
\setlength{\tabcolsep}{4pt}
\hspace*{-0.2cm}\begin{tabular}{@{}m{1.15cm}*{13}{>{\centering\arraybackslash}m{0.72cm}}@{}}
\toprule
& \rotatebox[origin=lb]{60}{\footnotesize AI / ML} & \rotatebox[origin=lb]{60}{\footnotesize Deep learning}
& \rotatebox[origin=lb]{60}{\footnotesize Forecasting} & \rotatebox[origin=lb]{60}{\footnotesize Anomaly / risk}
& \rotatebox[origin=lb]{60}{\footnotesize NLP / RAG} & \rotatebox[origin=lb]{60}{\footnotesize OCR / docs}
& \rotatebox[origin=lb]{60}{\footnotesize Blockchain} & \rotatebox[origin=lb]{60}{\footnotesize Escrow}
& \rotatebox[origin=lb]{60}{\footnotesize Marketplace} & \rotatebox[origin=lb]{60}{\footnotesize Ranking / MCDA}
& \rotatebox[origin=lb]{60}{\footnotesize Geo / GIS} & \rotatebox[origin=lb]{60}{\footnotesize Workflow} & \rotatebox[origin=lb]{60}{\footnotesize Dashboards} \\[1.9em]
\midrule
\rowcolor{gxteal!8}
\textbf{26132} & \yes & \pmid & \yes & \pmid & \pmid & \pmid & \yes & \yes & \yes & \yes & \pmid & \yes & \yes \\
\textbf{26100} & \yes & \non  & \non & \yes  & \yes  & \yes  & \yes & \pmid& \non & \pmid& \non & \yes & \yes \\
\textbf{26102} & \yes & \yes  & \pmid& \yes  & \pmid & \non  & \yes & \non & \non & \pmid& \pmid& \yes & \yes \\
\textbf{26021} & \yes & \pmid & \yes & \yes  & \non  & \pmid & \yes & \yes & \yes & \pmid& \non & \yes & \yes \\
\textbf{26002} & \yes & \yes  & \yes & \yes  & \non  & \pmid & \non & \non & \non & \yes & \yes & \yes & \yes \\
\textbf{26125} & \pmid& \non  & \non & \non  & \non  & \pmid & \yes & \yes & \non & \non & \non & \yes & \yes \\
\rowcolor{gxcyan!7}
\textbf{26130} & \yes & \non  & \non & \pmid  & \yes  & \yes  & \pmid& \non & \non & \pmid& \non & \yes & \yes \\
\rowcolor{gxcyan!7}
\textbf{26045} & \yes & \non  & \non & \non  & \yes  & \pmid & \non & \non & \non & \pmid& \non & \yes & \pmid \\
\textbf{26090} & \yes & \pmid & \yes & \non  & \yes  & \non  & \pmid& \pmid& \yes & \pmid& \non & \yes & \yes \\
\bottomrule
\end{tabular}

\vspace{1.2em}
\noindent\textbf{What each one still needs built}

\vspace{0.3em}
\renewcommand{\arraystretch}{1.35}
\noindent\begin{tabular}{@{}m{1.15cm}m{14.3cm}@{}}
\toprule
\textbf{PS} & \textbf{Net-new work} \\
\midrule
\rowcolor{gxteal!8}
26132 & \bld\ mandi price feed ingestion \quad \bld\ quality-grading capture \\
26100 & \emph{none} --- every stated requirement lands on existing capability \\
26102 & \emph{none} --- public scheme data is already published \\
26021 & \bld\ IoT hive sensor ingestion \quad \bld\ farmer mobile app \\
26002 & \bld\ live GPS telemetry \quad \bld\ offline sync \quad \bld\ multilingual notifications \\
26125 & \bld\ DID / verifiable credentials \quad \bld\ NFT minting \quad \bld\ on-chain RBAC \\
\rowcolor{gxcyan!7}
26130 & \bld\ department workflow connectors \\
\rowcolor{gxcyan!7}
26045 & \bld\ IP/GI corpus ingestion \quad \bld\ multilingual answering \\
26090 & \bld\ generative image enhancement \quad \bld\ voice-to-catalogue NLP \quad \bld\ native mobile app \\
\bottomrule
\end{tabular}

\clearpage
% ================= 26132 =================
\section{PS 26132 --- Market linkages and price discovery for farmers}
\psbanner{Government of Maharashtra (MSInS)}{Agriculture, FoodTech \& Rural Development}{Software}{\textbf{1 / 229}}{9 of 9}{20 Sep 2026}

\subsection*{Why this one}
Smallholders and producer groups cannot see current or expected prices across nearby mandis, processors
and institutional buyers. Information on quality specifications, demand, logistics, storage, payment
reliability and buyer credentials is scattered. Farmers sell immediately at harvest because they lack
storage and liquidity, and they negotiate from a weak position. Buyers, on the other side, cannot
aggregate consistent volumes or verify quality.

That is an information-asymmetry problem wrapped around a trust problem --- which is exactly the problem
GlobeX was built to solve, one market up. An exporter who cannot see which country will pay best, cannot
tell a reliable buyer from an unreliable one, and cannot be sure of getting paid, is the same person as a
farmer at a mandi gate. We already built: price and demand forecasting to answer \emph{when and where to
sell}, a multi-criteria ranking engine to turn many signals into one recommendation, counterparty matching
with trust scoring to answer \emph{who to sell to}, escrow with milestone conditions and dispute handling
to answer \emph{will I actually get paid}, and a tamper-evident ledger to make the record of it all
independently checkable.

The problem statement then asks, in its own words, for "transparent transaction records", "payment
reliability", "buyer credentials", "dispute or grievance processes" and "sale-window recommendations". We
have a working implementation of every one of those ideas. Nothing here has to be invented; it has to be
re-pointed.

\subsection*{How we would build it}
\renewcommand{\arraystretch}{1.35}
\noindent\begin{longtable}{@{}p{0.5cm}p{4.5cm}p{2.0cm}p{8.6cm}@{}}
\toprule
& \textbf{PS asks for} & \textbf{GlobeX} & \textbf{How it is built} \\
\midrule
\endhead
1 & Aggregate mandi prices, buyer demand, arrival volumes & Ingestion + feature pipeline & Point the existing trade-panel ingestion at Agmarknet/eNAM mandi feeds. Same shape of data: a price series per commodity per market per day, instead of per corridor per month. \\
2 & Localised price trends and \textbf{sale-window recommendations} & Forecasting (F1/F2) & The dual-head forecaster already predicts demand \emph{and} price together --- which is precisely what a "hold or sell now" recommendation needs. Use the quantile forecaster for the confidence band around the advice. \\
3 & Match farmers/FPOs with \textbf{verified} buyers & Matching + trust score & Semantic matching over commodity, grade, volume and location, then the composite trust score --- built for exactly this "is this counterparty credible" question. \\
4 & Quality specifications and grading & Document intelligence & Grade certificates and lot records go through the existing extraction and reconciliation path; grade claims are checked against the lot record rather than taken on trust. \\
5 & Lot creation and digital offers & Marketplace & Listings, offers and catalogue pages already exist; the seller becomes an FPO and the unit becomes a lot. \\
6 & Logistics and storage coordination & ETA + cost estimator & The distance/time and landed-cost model transfers directly; road transport constants replace ocean-freight ones. \\
7 & \textbf{Payment tracking} & Escrow lifecycle & Buyer funds escrow on offer acceptance; release is conditioned on delivery and grading, exactly like the existing DOCS / SHIPMENT / INSPECTION conditions. \\
8 & \textbf{Dispute or grievance process} & Dispute + resolve states & Already implemented as first-class states in the escrow contract, with the funds held rather than lost while a dispute is open. \\
9 & \textbf{Transparent transaction records} & Evidence ledger & Every offer, acceptance, grading and settlement is hash-anchored, so the record a farmer sees is the record the buyer sees and neither can quietly edit it. \\
\bottomrule
\end{longtable}

\begin{uspbox}{Our USP on 26132 --- escrow answers a problem the PS names out loud}
Most teams will build a price dashboard with a buyer directory. A price dashboard does not fix the
farmer's actual fear, which the problem statement states plainly: \textbf{payment reliability}. Our
differentiator is that the transaction is protected, not just informed --- the buyer's money is committed
into escrow before the lot moves, release is conditioned on delivery and grading, a dispute freezes the
funds instead of losing them, and the whole sequence is anchored so it can be audited later. That
converts "here is a better price signal" into "here is a trade you can safely complete", which is a
categorically stronger answer to the same problem statement. It is also the hardest part for another team
to build in a hackathon window --- and we already have it running.
\end{uspbox}

\begin{warnbox}{Risks to manage}
Mandi price data is public but messy --- inconsistent commodity naming across markets, missing days,
and units that vary by state; budget real time for the ingestion clean-up. Avoid over-claiming on the
forecaster: use the quantile model and present a range with a confidence band rather than a single
number, and never show a sale-window recommendation without the uncertainty around it.
\end{warnbox}

\clearpage
% ================= 26100 =================
\section{PS 26100 --- AI bid compliance verification for GeM procurement}
\psbanner{Ministry of Petroleum \& Natural Gas (CPCL)}{Smart Automation}{Software}{4 / 229}{14 of 14}{20 Sep 2026}

\subsection*{Why this one}
Procurement officers verify every bidder by hand across a dozen government registries --- Udyam, GSTN,
PAN, EPFO/ESIC, DigiLocker, debarment lists --- then decide whether the bid qualifies. It is slow,
inconsistent, and error-prone, and the consequence of getting it wrong is a contract awarded to an
ineligible vendor.

Strip the domain away and this is the problem GlobeX solves before every trade: take an entity, check it
against many authorities, reconcile the documents it submitted, score what came back, and give a human a
defensible recommendation with the evidence attached. All fourteen numbered requirements in the problem
statement land on something we have already built --- there is no net-new capability in it at all, which
is true of no other candidate except 26102.

\subsection*{How we would build it}
\renewcommand{\arraystretch}{1.35}
\noindent\begin{longtable}{@{}p{0.5cm}p{4.5cm}p{2.0cm}p{8.6cm}@{}}
\toprule
& \textbf{PS asks for} & \textbf{GlobeX} & \textbf{How it is built} \\
\midrule
\endhead
1 & Registry verification: Udyam, GSTN, PAN, EPFO, NSIC, Startup India & Entity screening & The screening engine already resolves an entity against a 31{,}629-record registry with fuzzy name matching; the registries swap, the matching logic does not. \\
2 & Blacklisting / debarment status & Denied-party screening & A debarment list \emph{is} a denied-party list. This is the same check, unmodified. \\
3 & Make in India / local content & Rules-of-origin logic & Local-content percentage thresholds are computationally identical to value-addition rules of origin, which the compliance store already encodes. \\
4 & Document verification and DigiLocker checks & OCR + reconciliation & Field extraction and cross-document consistency checking, already running over trade documents. \\
5 & Identify missing or inconsistent information & Reconciliation & The routine that catches an invoice quantity disagreeing with a bill of lading catches a turnover figure disagreeing with a filed return. \\
6 & Compliance score and bidder risk level & Composite scoring & Multi-factor scoring with banding, blending registry results and behavioural signals. \\
7 & AI recommendation, officer keeps the decision & Transaction gate + dossier & A human-in-the-loop decision gate with a machine recommendation and a written justification --- already how trade approval works. \\
8 & \textbf{Auditable record of every check} & Evidence ledger & Each verification is hash-anchored with its result and timestamp. \\
\bottomrule
\end{longtable}

\begin{uspbox}{Our USP on 26100 --- an audit trail that cannot be edited after the fact}
Requirement 14 asks for an auditable record. Most teams will satisfy that with a database log --- which
the awarding authority can edit. Ours is cryptographically anchored, so a tender decision made in March
can be proved in November to have been based on exactly the evidence shown at the time. In public
procurement, where the dispute is usually \emph{"what did you actually check before you awarded this?"},
that is the difference between a log and proof. This is the one candidate where blockchain is demanded by
the problem statement rather than added by us, so it cannot be dismissed as decoration.
\end{uspbox}

\begin{warnbox}{Risks to manage}
Live portal credentials (GSTN, DigiLocker) will not be issued to a student team --- demo against the dummy
bidder and tender datasets the problem statement itself offers, and draw the integration boundary
explicitly in the architecture slide rather than implying live access.
\end{warnbox}

\clearpage
% ================= 26102 =================
\section{PS 26102 --- Anomaly, fraud and inefficiency detection in MPLADS}
\psbanner{MoSPI --- Data Informatics \& Innovation Division}{Smart Automation}{Software}{31 / 229}{9 of 9}{20 Sep 2026}

\subsection*{Why this one}
Thousands of local works are funded through many implementing agencies, producing large volumes of
expenditure and execution data that nobody can review manually. The problem statement asks for ML that
surfaces anomalies in expenditure, fund utilisation, cost estimates and work execution, so irregularity is
caught early.

This is the closest match in the dataset to our strongest asset. The trade anomaly classifier is trained,
chronologically validated and performing at 0.98 F1, and the unsupervised risk ensemble exists precisely
because trade data has \emph{no labelled fraud} --- the same condition public-fund data is in. Most teams
attacking this PS will have to invent an approach for unlabelled irregularity detection during the
hackathon. We already have one, with attribution behind every flag.

\subsection*{How we would build it}
\renewcommand{\arraystretch}{1.35}
\noindent\begin{longtable}{@{}p{0.5cm}p{4.5cm}p{2.0cm}p{8.6cm}@{}}
\toprule
& \textbf{PS asks for} & \textbf{GlobeX} & \textbf{How it is built} \\
\midrule
\endhead
1 & Anomalies in expenditure patterns & XGBoost anomaly + features & Sanctioned-work expenditure records are the same object as trade transactions: an amount, a date, a category, a counterparty. The rolling-statistics and growth features transfer directly. \\
2 & Irregularity without labelled fraud & Isolation Forest + autoencoder & The unsupervised ensemble runs unchanged. Reconstruction error over a work's spending sequence flags what does not resemble normal execution. \\
3 & Cost-estimate deviation & Unit-value deviation & Our peer-price deviation feature becomes cost-per-unit against comparable works in comparable districts. \\
4 & Duplicate or overlapping works & TF-IDF similarity & Near-duplicate work descriptions across constituencies are caught the same way near-equivalent company records are. \\
5 & Risk scoring and prioritisation & Composite score & Ranks works and agencies so limited audit capacity goes where the signal is strongest. \\
6 & Explainable flags & TreeSHAP attribution & Every flag carries its contributing factors --- non-negotiable when the subject is public money. \\
7 & Monitoring dashboard and alerts & Dashboards + n8n & Existing charting stack and resilient workflow orchestration. \\
\bottomrule
\end{longtable}

\begin{uspbox}{Our USP on 26102 --- blockchain as the accountability layer, not decoration}
Here is the move most teams will miss. An anomaly-detection tool that accuses an implementing agency
creates a new problem: who flagged this, on what evidence, and did anyone quietly change it afterwards?
We anchor \textbf{the flag itself} --- the model version, the input snapshot hash, the score and the
officer's disposition --- to the evidence ledger. That makes the oversight trail as tamper-evident as the
spending it oversees, protects agencies from silently altered accusations, and protects auditors from
claims of tampering.

\vspace{0.5em}
Paired with our second differentiator --- we do \textbf{not} claim the model detects fraud. Our own risk
documentation refuses to equate an anomaly with fraud because there is no ground truth for it, so the
output is "this pattern is unusual and warrants review", feeding a human investigation. In a domain where
a false public accusation is a real harm, that restraint plus an immutable audit trail is a far stronger
pitch than a higher accuracy number.
\end{uspbox}

\begin{warnbox}{Risks to manage}
Do not present the 0.98 F1 from trade data as a result on MPLADS data --- it is evidence the method works,
not a claim about this domain, and a judge who catches that conflation will discount everything else.
Retrain and report honestly on the scheme's own published data.
\end{warnbox}

\clearpage
% ================= 26021 =================
\section{PS 26021 --- Honey Chain: blockchain traceability and smart beekeeping}
\psbanner{Ministry of MSME (KVIC Honey Mission)}{Agriculture, FoodTech \& Rural Development}{Software}{30 / 229}{7 of 8}{20 Sep 2026}

\subsection*{Why this one}
KVIC-supported beekeepers face counterfeit honey diluting their market, weak buyer linkage, and no way to
prove their product is what they say it is. The problem statement asks for blockchain-based batch
traceability with QR consumer verification, plus AI/IoT support for disease detection and yield
prediction.

Structurally this is GlobeX's spine with a different commodity: producers listed in a marketplace,
provenance anchored on-chain, forecasting on the supply side, anomaly detection on the operational side.
The chain layer here is not bolted on --- authenticity \emph{is} the problem, which makes it the most
natural home for our blockchain work of any candidate except 26100.

\subsection*{How we would build it}
\renewcommand{\arraystretch}{1.35}
\noindent\begin{longtable}{@{}p{0.5cm}p{4.5cm}p{2.0cm}p{8.6cm}@{}}
\toprule
& \textbf{PS asks for} & \textbf{GlobeX} & \textbf{How it is built} \\
\midrule
\endhead
1 & Batch traceability, hive to shelf & Chain anchor + ledger & Each batch event (harvest, extraction, packing, dispatch) is hash-anchored using the existing write-intent-then-confirm pattern --- the payload changes, the guarantee does not. \\
2 & QR consumer verification & Hash re-verification & Scanning re-computes the hash and compares it against the anchored record: our document integrity check, exposed through a QR code. \\
3 & Market linkage for producers & Marketplace + matching & Listings and buyer matching already exist; sellers become beekeeper collectives. \\
4 & Productivity / yield prediction & Quantile forecaster & Per-cluster yield forecasting with confidence bands. Use the quantile model, not the GRU. \\
5 & Disease and environmental monitoring & Anomaly detection & Anomaly detection over hive sensor series is the same problem as over trade series; only the ingestion path is new. \\
6 & Authenticity and market credibility & Escrow & See USP below. \\
\bottomrule
\end{longtable}

\begin{uspbox}{Our USP on 26021 --- payment released against verified quality}
"Blockchain + AI + QR for traceability" is the most crowded pattern in any hackathon; assume several
teams present it. What they will not have is the second half of the trust problem. Traceability proves
where honey came from but does nothing about whether the beekeeper actually gets paid fairly for it. We
add \textbf{quality-conditioned escrow}: the buyer commits funds when the batch is listed, release is
conditioned on the inspection result recorded against that batch, and a dispute holds the money rather
than leaving the producer to chase it. Provenance plus guaranteed settlement is a complete answer to
"why does counterfeit honey persist" --- provenance alone is only half of it.
\end{uspbox}

\begin{warnbox}{Risks to manage}
There is no hive sensor dataset available --- the IoT layer must be demonstrated on simulated or public
apiary data, stated openly. And because the pattern is crowded, the pitch must lead with root cause (why
existing certification has failed to stop dilution) rather than with the technology stack.
\end{warnbox}

\clearpage
% ================= 26125 =================
\section{PS 26125 --- Blockchain identity, access control and digital asset management}
\psbanner{Bharat Electronics Limited}{Blockchain \& Cybersecurity}{Software}{\textbf{1 / 229} on the chain profile}{6 of 8}{20 Sep 2026}

\subsection*{Why this one}
Centralised identity and access management is a single point of failure: vulnerable to breach, identity
theft and unauthorised access, and it leaves asset ownership scattered across disconnected systems where
authenticity and ownership history cannot be reliably verified. The problem statement asks for a
decentralised alternative --- self-sovereign identifiers authenticated by cryptographic proof, digital
assets issued as NFTs bound permanently to those identities, smart contracts enforcing who may mint,
allocate and transfer, role-based access control enforced on-chain rather than in application code, and
every operation immutably recorded as an audit trail.

This is the deepest chain-layer match in the dataset --- first of 229 on the blockchain profile, roughly
70\% clear of second place. The reason is that GlobeX's escrow contract already embodies the exact
discipline this statement is asking for: state transitions governed by a contract rather than by
application code, actions permitted only to authorised parties, and a rule that no consequential state is
ever written without a confirmed on-chain receipt behind it. That last point is the difference between a
system that \emph{uses} a blockchain and one that can actually be \emph{trusted} because of it, and we
already enforce it in production code.

\vspace{0.4em}
\noindent The honest counterweight: this problem statement contains \textbf{almost no machine learning}
--- it ranks 200th of 229 on the ML profile. Choosing it means our forecasting, anomaly detection, risk
ensemble and ranking work --- the majority of the engineering behind GlobeX --- does not appear in the
solution at all. It is the exact mirror image of PS 26102's trade-off.

\subsection*{How we would build it}
\renewcommand{\arraystretch}{1.35}
\noindent\begin{longtable}{@{}p{0.5cm}p{4.5cm}p{2.0cm}p{8.6cm}@{}}
\toprule
& \textbf{PS asks for} & \textbf{GlobeX} & \textbf{How it is built} \\
\midrule
\endhead
1 & Smart contracts governing all operations & Escrow contract + adapter & We already run a lifecycle contract with guarded state transitions through a dedicated chain-adapter service; the contract's subject changes from a trade to an identity and its assets. \\
2 & Immutable audit trail of every operation & Evidence ledger & Identity creation, minting, allocation, transfer and permission changes are anchored exactly as document and trade events are anchored today. \\
3 & Ownership and authenticity verification & Hash verification & Re-deriving a hash and comparing it against the anchored record is our existing integrity check. \\
4 & Controlled issuance --- only authorised admins may mint & Role-gated operations & The escrow contract already restricts who may fund, release and resolve; the same authorisation pattern governs minting. \\
5 & Ownership history & Ledger query + UI & The blockchain ledger screen already renders anchored event history for a trade; here it renders it for an asset. \\
6 & Transaction confirmation discipline & Write-intent pattern & Our submit-then-confirm rule (never record a final state on an unconfirmed transaction) transfers wholesale, and is the single most valuable thing we bring. \\
7 & Decentralised identifiers \& verifiable credentials & \bld\ new & DID issuance and cryptographic proof verification is new contract and client work. \\
8 & NFT minting and on-chain RBAC & \bld\ new & Token standard implementation and on-chain role enforcement go beyond our current escrow contract. \\
\bottomrule
\end{longtable}

\vspace{-0.4em}
\noindent{\small\textbf{Coverage: 6 of 8 asks land on existing capability}; 2 require new Solidity work.}

\begin{uspbox}{Our USP on 26125 --- the discipline, not the buzzwords}
Every team attempting this problem statement will demo a contract that mints an NFT. Most will also,
somewhere in the stack, write "transferred" into a database before the chain has confirmed it --- because
that is the shortcut that makes a demo work on a flaky testnet. Our differentiator is the failure
handling nobody else will have thought about: we already write a pending intent record \emph{before} the
chain call, patch it to confirmed with the real transaction hash on success or to failed with the real
error on failure, and never let the application claim an ownership change the chain has not
acknowledged. In a problem statement whose entire premise is trustworthiness, being the one team whose
system cannot lie about its own state is a far stronger position than having the fanciest token design.
\end{uspbox}

\begin{warnbox}{The trade-off to weigh}
This is the narrowest of the nine candidates. It exercises our chain layer more deeply than anything
else, and nothing else at all --- no forecasting, no anomaly detection, no risk scoring, no ranking, no
compliance retrieval. If the team's judgement is that GlobeX's identity is the blockchain work, this is
the strongest possible showcase for it. If the ML work is the identity, this is the weakest of the nine.
\end{warnbox}

\clearpage
% ================= 26002 =================
\section{PS 26002 --- AI smart logistics and accessibility platform (NER)}
\psbanner{Ministry of Development of North Eastern Region}{Transportation \& Logistics}{Software}{40 / 229}{6 of 8}{20 Sep 2026}

\subsection*{Why this one}
Landslides, floods and fragile terrain cut road links to remote NER districts, delaying medicines, food
and produce. The problem statement asks for a platform that monitors accessibility, predicts disruptions,
recommends alternate routes, tracks vehicles, alerts authorities, accepts geo-tagged field reports and
presents district-level dashboards.

The analytical core maps well, and one mapping is genuinely elegant: \textbf{ranking alternate routes is
the same problem as ranking destination markets}. Our multi-criteria engine already turns eight weighted,
explainable signals into one ranked recommendation --- point it at route options scored on rainfall risk,
slope, incident history, detour cost and time, and it does the same job. Corridor-level risk scoring is
also already implemented; today the corridors are trade lanes, here they are road segments.

\subsection*{How we would build it}
\renewcommand{\arraystretch}{1.35}
\noindent\begin{longtable}{@{}p{0.5cm}p{4.5cm}p{2.0cm}p{8.6cm}@{}}
\toprule
& \textbf{PS asks for} & \textbf{GlobeX} & \textbf{How it is built} \\
\midrule
\endhead
1 & Predict route disruptions & Forecast + anomaly + risk & The same three-model shape we run on trade corridors, with rainfall, slope and incident history replacing trade features. \\
2 & Alternate route ranking + travel delay & MCDA ranking + ETA & Weighted, explainable route scoring; delay estimated by the distance-over-speed method with terrain and condition penalties. \\
3 & High-risk corridor identification & Risk ensemble & Corridor-level risk scoring, already implemented --- the object the PS calls a "high-risk transport corridor" is the object K1 already produces. \\
4 & Automated alerts & n8n workflows & Resilient multi-step alerting with per-step status, already built. \\
5 & District connectivity dashboards & Dashboards + geo view & Existing charting and geospatial visualisation stack. \\
6 & Geo-tagged field reports & Document intake & Upload, extraction and persistence exist; geo-tagging and image handling are additions. \\
7 & GPS vehicle tracking & \bld\ new & Live telemetry ingestion, position storage and streaming. Our ETA service is a static estimator, not tracking. \\
8 & Offline sync + multilingual alerts & \bld\ new & Offline-first synchronisation for low-network areas is a discipline we have not built. \\
\bottomrule
\end{longtable}

\begin{uspbox}{Our USP on 26002 --- explainable route decisions, not a black box}
Most teams will produce a map with a shortest-path line on it. Ours ranks routes on weighted criteria and
\emph{shows why} --- this route scores lower because of slope-adjusted landslide risk and recent incident
density, not merely because it is longer. For a district officer deciding whether to send a medicine
convoy, an explainable, auditable recommendation is the difference between a tool they trust and one they
override. Our backtested ranking methodology is the credible core here.
\end{uspbox}

\begin{warnbox}{Be honest about two things}
Live GPS tracking is a headline feature of this platform and we would be building it from zero --- say so
rather than implying we have it. And the chain and escrow half of GlobeX has no natural role here; forcing
it in reads as technology for its own sake, which the first-round rubric penalises. This is the one
candidate where our blockchain work is a spectator.
\end{warnbox}

\clearpage
% ================= 26090 =================
\section{PS 26090 --- AI market linkage and smart cataloging for artisans}
\psbanner{Ministry of Social Justice and Empowerment}{Heritage \& Culture}{Software}{15 / 229}{6 of 9}{20 Sep 2026}

\subsection*{Why this one}
Marginalised artisans and weavers get a sales bump from physical exhibitions but have no year-round
digital channel, blocked by low digital literacy, language barriers and the practical difficulty of
photographing, pricing and cataloguing products to e-commerce standards. The ask is a mobile app acting as
a "virtual business manager": AI image enhancement, voice-to-catalogue in regional languages, and a
dynamic pricing assistant, connecting artisans to B2B buyers and government e-marketplaces.

The commercial half of this is our marketplace: catalogue and listing creation, buyer discovery and
matching, and price recommendation --- all of which exist, including the listing and catalogue editor
screens. The reason it sits below the others is the other half: generative image enhancement,
speech-to-text in regional languages and generated product copy are capabilities we have not built, and
the deliverable is explicitly a native mobile app rather than a web platform.

\subsection*{How we would build it}
\renewcommand{\arraystretch}{1.35}
\noindent\begin{longtable}{@{}p{0.5cm}p{4.5cm}p{2.0cm}p{8.6cm}@{}}
\toprule
& \textbf{PS asks for} & \textbf{GlobeX} & \textbf{How it is built} \\
\midrule
\endhead
1 & Dynamic pricing assistant & Forecasting + ranking & Price recommendation from comparable-product signals with a confidence band --- our forecasting and scoring path, applied to craft categories. \\
2 & Connect artisans to B2B buyers / GeM & Marketplace + matching & Listings, buyer discovery, semantic matching and trust scoring already exist. \\
3 & Product cataloguing and inventory & Catalogue editor & Catalogue and listing creation screens are already built. \\
4 & Buyer credibility for vulnerable sellers & Trust scoring + escrow & See USP below. \\
5 & AI image enhancement studio & \bld\ new & Background removal and lighting correction --- a generative/CV capability we do not have. \\
6 & Voice-to-catalogue in regional languages & \bld\ new & Speech recognition plus generated multilingual product copy. \\
7 & Native cross-platform mobile app & \bld\ new & Our frontend is web; this PS explicitly wants mobile-first for low-literacy users. \\
\bottomrule
\end{longtable}

\begin{uspbox}{Our USP on 26090 --- protecting the seller, not just listing them}
Every team will build the camera and the auto-description; those features are the obvious reading of this
problem statement. The risk nobody addresses is that a low-literacy artisan, newly exposed to B2B buyers,
is the party most vulnerable to non-payment and exploitative terms. Our differentiator is
\textbf{verified buyers plus protected payment} --- trust-scored buyer credentials and escrowed settlement
so the artisan ships against committed funds rather than a promise. For a Ministry of Social Justice
problem statement, protecting the beneficiary from the market they are being introduced to is a
substantially more thoughtful answer than a prettier product photo.
\end{uspbox}

\begin{warnbox}{Why it is ranked sixth}
Three of the nine asks are net-new capability, and two of them (generative imaging, multilingual speech)
are the features the problem statement leads with. We would be strongest on the half of the problem the
judges may treat as secondary --- a real risk of a mismatch between our strength and the rubric's weight.
\end{warnbox}

\clearpage
% ================= NEW FINDS =================
\section{New finds from the extended sweep}
\textit{Neither of these appeared in the original product-shaped search. Both were surfaced by the
compliance/RAG profile, and both are strong enough to sit in the main comparison. The problem-statement
Excel was checked row by row against the searched dataset first --- 229 IDs, no differences in title,
description, category or theme --- so this sweep covers the complete published set.}

\subsection*{PS 26130 --- Streamlining industrial approvals, compliance and support services}
\psbanner{Government of Maharashtra (MSInS)}{Miscellaneous}{Software}{9 / 229 on the compliance profile}{8 of 8}{20 Sep 2026}

\noindent\textbf{Why.} Entrepreneurs setting up an industrial unit must obtain registrations, licences,
NOCs, inspections and renewals from many authorities, and the requirements vary by sector, location,
project size and stage of operation. Applicants cannot work out which approvals apply to them; departments
receive incomplete applications and cannot see where the bottlenecks are.

This is PS 26100 seen from the other side of the counter. Where 26100 asks us to verify a bidder against
authorities, 26130 asks us to tell an applicant which authorities apply and walk them through it --- and
the core of that request, \emph{"generate a customised approval checklist"}, is precisely what our
required-document matrix does today: take an item's classification plus its origin and destination, and
return the exact set of documents and clearances that apply to it. Swap HS6 code and trade lane for sector
and location and the function is unchanged.

\vspace{0.5em}
\noindent\textbf{How.} The applicable-approval checklist comes from the document matrix; documentation
guidance comes from the citation-backed retrieval engine; pre-validation of submissions reuses the
extraction and inconsistency-detection path; "reuse verified data" is our verified-entity record; parallel
departmental coordination, inspection scheduling and SLA timers are n8n workflows with per-step status
already built for exactly that failure-isolation reason; bottleneck visibility is the dashboard layer.
All eight stated outcomes land on existing capability --- only the department-side connectors are new.

\begin{uspbox}{USP on 26130 --- the checklist is the hard part, and we already built it}
Most teams will build a portal: forms, uploads, a status tracker. The genuinely difficult piece is the one
the problem statement leads with --- correctly deriving \emph{which} approvals apply to a given unit, when
the answer depends on four interacting variables. That is a rules-and-retrieval problem we have already
solved once for trade compliance. Adding chain-anchored approval records on top means an issued clearance
cannot be quietly back-dated or altered, which matters in exactly the environment where approvals are
contested.
\end{uspbox}

\subsection*{PS 26045 --- IP-SAKTI Sahayak: multilingual, source-cited RAG assistant}
\psbanner{Ministry of Ayush}{MedTech / BioTech / HealthTech}{Software}{4 / 229 on the compliance profile}{6 of 7}{20 Sep 2026}

\noindent\textbf{Why.} Protecting and commercialising an Ayurvedic product means navigating several
overlapping legal regimes simultaneously --- patents, geographical indications, trademarks, copyright,
designs, trade secrets, plant-variety rights, access-and-benefit-sharing duties over biological resources,
and a drug-regulatory framework that decides whether a formulation counts as a classical medicine, a
proprietary medicine, a new drug, a phytopharmaceutical, a food or a cosmetic. The ask is a multilingual
assistant that answers in plain language \textbf{with citations to the specific source}.

That is a closer description of our RAG engine's actual job than any other statement in the dataset,
including PS 26107. Our retriever already spans multiple overlapping regulatory corpora at once --- tariff
schedules, a sanctions registry, export-control lists, rules of origin, SPS/TBT standards --- and answers
against them with provenance attached. The classification question is familiar too: deciding whether a
formulation is a classical medicine or a phytopharmaceutical, and therefore which pathway applies, is the
same shape as deciding an HS6 code and therefore which document set applies.

\vspace{0.5em}
\noindent\textbf{How.} Retrieval, ranking and citation enforcement transfer directly. The versioned
current-facts store matters more here than anywhere else, because this domain moved recently --- the 2024
patent and biodiversity rules and the 2024 WIPO treaty --- and an assistant citing superseded law is worse
than no assistant. New work: ingesting the IP, GI and AYUSH regulatory corpora, and answering in multiple
languages.

\begin{uspbox}{USP on 26045 --- we already refuse to answer without a source}
The obvious build here is a chatbot over some documents, and the obvious failure is a confident answer
with an invented citation --- which, in a legal and regulatory domain, is not a glitch but a liability for
the person who acted on it. Our engine is built the other way round: provenance is mandatory and
fabricated citations are prohibited by design. Leading with \emph{refusal behaviour} --- what the
assistant does when it does not have a source, and how it shows which clause an answer rests on --- is a
credibility argument almost no competing team will make, and it is the argument that matters most to a
ministry putting its name on legal guidance.
\end{uspbox}

\clearpage
% ================= DECISION =================
\section{Recommendation}

\renewcommand{\arraystretch}{1.5}
\noindent\begin{tabular}{@{}p{1.2cm}p{3.6cm}p{2.3cm}p{2.6cm}p{5.6cm}@{}}
\toprule
\textbf{PS} & \textbf{Best when the team wants to lead with} & \textbf{Net-new work} & \textbf{Blockchain role} & \textbf{Headline USP} \\
\midrule
\rowcolor{gxteal!10}
26132 & The whole product & Low & Central & Escrow answers "payment reliability", named in the PS \\
26100 & Compliance \& verification & \textbf{None} & Demanded by PS & Audit trail that cannot be edited afterwards \\
26102 & The trained ML core & \textbf{None} & Added by us & Tamper-evident record of the flags themselves \\
26021 & Blockchain \& traceability & Medium & Central & Payment released against verified quality \\
26002 & Analytics \& geospatial & \textbf{High} & \cellcolor{gxamber!12}None honest & Explainable route ranking, not a black box \\
26125 & Blockchain, purely & Medium & \cellcolor{gxviolet!12}Everything & Escrow-grade contract discipline: no state written without a confirmed receipt \\
\rowcolor{gxcyan!7}
26130 & Compliance, applicant side & Low & Optional & Approval checklist generation --- our required-document matrix, re-pointed \\
\rowcolor{gxcyan!7}
26045 & RAG with citations & Medium & None & Refusing to answer without a source --- already enforced in our engine \\
26090 & Marketplace \& social impact & \textbf{High} & Light & Verified buyers + protected payment for artisans \\
\bottomrule
\end{tabular}

\vspace{1.4em}
\begin{callout}{The recommendation}
\textbf{Lead with PS 26132.} It is the measured closest match to GlobeX in the entire dataset by a wide
margin, every one of its nine asks lands on something we have built, and the differentiator is not a
technology we are showing off but the direct answer to a problem the statement itself raises --- farmers
selling into a market where payment is not reliable. It also lets us present the full platform story
(intelligence $\rightarrow$ matching $\rightarrow$ transaction $\rightarrow$ settlement $\rightarrow$
dispute) as one coherent solution rather than a collection of features.

\vspace{0.6em}
\textbf{Keep PS 26100 as the strong second.} It is the only other candidate with zero net-new work, and
it is the one where blockchain is required by the problem statement rather than offered by us --- the
safest possible ground for defending that part of the stack under questioning.

\vspace{0.6em}
\textbf{PS 26102 is the pick if the team would rather lead with the models}, and the blockchain angle
(anchoring the flags themselves) turns what looked like an unused half of GlobeX into that submission's
sharpest idea.

\vspace{0.6em}
\textbf{PS 26002 remains viable} and the route-ranking insight is real, but it carries the most net-new
work of the six and leaves our most distinctive asset on the bench.
\end{callout}

\vspace{0.9em}
\noindent\textbf{Where the three additions land.} \textbf{PS 26125} is the strongest showcase for the
chain layer specifically --- first of 229 on that profile --- but it is the narrowest option on the table,
using none of the ML work. \textbf{PS 26130} is the most interesting of the two new finds: it matches
26100's low-drift profile from the applicant's side, all eight of its outcomes land on existing
capability, and its hardest requirement is a function we have already built. \textbf{PS 26045} is the
truest match to the RAG engine and displaces 26107 as the retrieval-led option, though its corpus
ingestion is real work. None of the three displaces 26132 at the top.

\vspace{0.9em}
\noindent\textbf{One practical note.} 26132, 26021 and 26090 all exercise the same marketplace, matching,
escrow and ledger path, so preparation for any one transfers substantially to the others. 26100, 26130 and
26045 share the screening, retrieval, document and evidence path --- that cluster is now three statements
deep, which makes it the safest area to invest preparation in. 26102 shares the scoring half of it. The
two outliers are 26002 (its telemetry work transfers nowhere) and 26125 (its contract work transfers only
to 26021), so both should be decided on early rather than prepared in parallel.

\vspace{0.9em}
\noindent\textbf{A note on the search itself.} The problem-statement Excel supplied by the team was
verified against the dataset used here before this sweep was run: 229 rows, 229 unique IDs, and no
difference in any title, description, category or theme. Every candidate in this report was therefore
drawn from the complete published set, scored under four independent capability profiles.

\vspace{1.2em}
\begin{warnbox}{Unrelated to PS choice, but worth fixing before any submission}
The repository's new \path{README.md} advertises a demand-forecast error of \textbf{MAPE 8.42\%}
("Deep GRU + Quantile Loss"). The benchmark artifact still sitting in both model folders,
\path{benchmark_comparison.csv}, reports the Dual-Head GRU at \textbf{54.49\% MAPE, R² 0.399} --- worse
than a three-year moving average on the same data. Either a newer quantile model exists and the benchmark
file is stale, or the headline figure is not backed by the artifact. Whichever it is, the two should be
reconciled before that number appears in a submission, because it is the first thing a technical judge
would check.
\end{warnbox}

\vspace{1.5em}
{\footnotesize\color{gxslate}\textit{Candidates selected by TF-IDF cosine similarity of a GlobeX
capability profile against all 229 problem statements from \url{https://www.sih.gov.in/sih2026PS}, then
mapped against verified backend modules, trained model artifacts and frontend screens. GlobeX is an
in-flight build --- this maps capability direction, not final code. All six share a 20 September 2026
idea-submission deadline.}}

\end{document}
